Դիտողություն ($\mathbb{R}^\infty$ տարածության մասին)$\mathbb{R}^n$ էվկլիդյան տարածությունների փակ ենթատարածությունների օրինակ 6-ում բերված հատկությունը թույլ է տալիս սահմանել $\mathbb{R}^\infty$ բազմությունը և նրա վրա հատուկ տոպոլոգիա: Սահմանենք $\mathbb{R}^\infty$ բազմությունը որպես $\mathbb{R}^1 \cup \mathbb{R}^2 \cup \mathbb{R}^3 \cup ...$ բազմության ֆակտոր-բազմություն կատարելով նրա տարրերի$$(x_1, x_2, ..., x_n) = (x_1, x_2, ..., x_n, 0) = (x_1, x_2, ..., x_n, 0, 0) = ...$$տեսքի նույնացումներ ամեն մի բնական $n$ թվի և իրական թվերի ամեն մի $(x_1, x_2, ..., x_n)$ հաջորդականության համար: Թույլ տալով մեզ որոշ ազատություն կարող ենք համարել որ $\mathbb{R}^\infty$-ը իրական թվերից կազմված բոլոր անվերջ $(x_1, x_2, ...)$ հաջորդականությունների բազմությունն է, որոնց անդամները սկսած որոշ տեղից զրոներ են: Միանգամայն այս եղանակով $\mathbb{R}^\infty$-ը վերածվում է անվերջ չափանի իրական գծային տարածության հաջորդականությունների անդամ առ անդամ գումարման և թվով բազմապատկման գործողություններով: Սահմանենք տոպոլոգիա $\mathbb{R}^\infty$-ի վրա համարելով փակ ամեն մի $F \subset \mathbb{R}^\infty$ ենթաբազմություն այն և միայն այն դեպքում, երբ ամեն մի $n$-ի դեպքում փակ է $F \cap \mathbb{R}^n$ հատումը: Տոպոլոգիայի F1-F3 աքսիոմները ստուգվում են հեշտությամբ:$\mathbb{R}^\infty$-ում կարելի է սահմանել տոպոլոգիա ուրիշ եղանակով՝ վերածելով $\mathbb{R}^\infty$-ը մետրիկական տարածության $(\rho(x, y)^2 = \sum_{i \ge 1} (x_i - y_i)^2$ մետրիկով):Նշենք, որ $\mathbb{R}^\infty$-ի այս երկու տոպոլոգիաները համարժեք չեն, քանի որ, օրինակ$$A = \{(1, 0, 0, ...), (0, \frac{1}{2}, 0, 0, ...), (0, 0, \frac{1}{3}, 0, 0, ...), ...\}$$ենթաբազմությունը ակնհայտորեն փակ է $\mathbb{R}^\infty$-ի առաջին տոպոլոգիայում, բայց փակ չէ երկրորդ տոպոլոգիայում: Իրոք, $O=(0, 0, 0, ...)$ կետի ցանկացած $\varepsilon$-շրջակայք ունի ոչ դատարկ հատում $A$-ի հետ, բայց $O$-ն չի պատկանում $A$-ին (հիմնավորեք):Սահմանվում է $\mathbb{R}^\infty$ բազմություն և նրա վրա տոպոլոգիա ևս մի եղանակով. այս դեպքում $\mathbb{R}^\infty$-ի տարրեր են համարվում բոլոր $(x_1, x_2, ...)$ տեսքի հաջորդականությունները, պայմանով որ զուգամիտում է $\sum_{i=1}^\infty x_i^2$ շարքը: Այս դեպքում նույնպես $\mathbb{R}^\infty$-ը վերածվում է մետրիկական տարածության, որը նշանակվում է $l_2$-ով և կոչվում է հիլբերտյան տարածություն (մանրամասները տես [ ]-ում):Այժմ ներկայացնենք մակածված տոպոլոգիա և ենթատարածություն հասկացություններին առնչվող մի քանի պարզ և հաճախ կիրառվող փաստեր:Խնդրի շարունակում (Extension Problem)(Red text note)Հաճախ դիտարկվում է նաև հակառակ խնդիրը. տրված $g: A \to Y$ անընդհատ արտապատկերումը շարունակել մինչև $f: X \to Y$ անընդհատ արտապատկերում: (Իհարկե լուծում չունի միշտ):Նկատենք, որ ընդհանուր դեպքում անընդհատ շարունակում կարող է գոյություն չունենալ:Տոպոլոգիական դասակարգում (Topological Classification)Վերը բերված 1-6 օրինակները ցույց են տալիս, որ տոպոլոգիական տարածություն ենթատարածություն հասկացությունը հնարավորություն է տալիս արդեն հայտնի տարածություններից ստանալ բազմաթիվ նոր տարածություններ դրանց ենթաբազմությունների տեսքով:Վերլուծական երկրաչափություն դասընթացում կատարվում է 2-րդ կարգի կորերի և 2-րդ կարգի մակերևույթների մետրիկական և աֆինական դասակարգումներ:Յուրաքանչյուր 2-րդ կարգի կոր կարող ենք դիտել որպես տոպոլոգիական տարածություն, հարթության սովորական մետրիկական տոպոլոգիայից կորի վրա մակածված տոպոլոգիայով այնպես, ինչպես դա արված է օրինակ 3-ում շրջանագծի դեպքում: Իսկ յուրաքանչյուր 2-րդ կարգի մակերևույթ կարող ենք դիտել որպես տոպոլոգիական տարածություն $\mathbb{R}^3$ էվկլիդյան տարածության տոպոլոգիայից մակերևույթի վրա մակածված տոպոլոգիայով այնպես, ինչպես դա նկարագրված է օրինակ 4-ում սֆերայի դեպքում:Առաջանում է խնդիր. կատարել 2-րդ կարգի կորերի և մակերևույթների տոպոլոգիական դասակարգումներ:Քանի որ իզոմետրիկ և աֆինական ձևափոխությունները փոխմիարժեք և անընդհատ արտապատկերումներ են, ուստի միմյանց մետրիկորեն կամ աֆինորեն համարժեք ցանկացած երկու 2-րդ կարգի կոր կամ մակերևույթ նաև տոպոլոգիորեն համարժեք են (հոմեոմորֆ են):Նկատենք, որ աֆինորեն ոչ համարժեք երկու կոր կամ մակերևույթ կարող են լինել տոպոլոգիորեն համարժեք: Օրինակ, ցանկացած հիպերբոլ հոմեոմորֆ է զուգահեռ իրական ուղիղների զույգի:Իրոք, տեղադրելով հիպերբոլն ու ուղիղների զույգը միմյանց նկատմամբ հարմար դիրքով, կարող ենք զուգահեռ պրոյեկտել դրանցից մեկը մյուսի վրա ինչպես ցույց է տրված ստորև գծագրում:Կորերի և մակերևույթների լիարժեք տոպոլոգիական դասակարգումը ենթադրում է ոչ միայն հաստատել, որ տվյալ երկու կորը կամ մակերևույթը հոմեոմորֆ են, այլ նաև ցանկացած երկու կորի կամ մակերևույթի դեպքում հիմնավորել նրանց ոչ հոմեոմորֆությունը, եթե չի հաջողվում կառուցել հոմեոմորֆիզմ նրանց միջև:Օրինակ, ինչպես գիտենք, որևէ էլիպս աֆինորեն համարժեք չէ որևէ հիպերբոլի, քանի որ աֆինական ձևափոխությունները սահմանափակ պատկերները արտապատկերում են սահմանափակ պատկերների (էլիպսը սահմանափակ պատկեր է, հիպերբոլը՝ ոչ):Աֆինական ձևափոխությունների այս հատկությունը կարելի է վերաձևակերպել նաև այսպես՝ պատկերի սահմանափակությունը (ինչպես նաև անսահմանափակությունը) աֆինական հատկություն է:Բայց (ինչպես դա հետևում է ստորև 12.6 խնդրից) պատկերի սահմանափակությունը չի հանդիսանում տոպոլոգիական հատկություն:Ուստի հիմնավորելու համար, որ էլիպսը և հիպերբոլը տոպոլոգիապես համարժեք չեն (հոմեոմորֆ չեն) անհրաժեշտ է գտնել պատկերների որևէ տոպոլոգիական հատկություն, որով օժտված լինի այդ պատկերներից միայն մեկը: Այդպիսի որոշ տոպոլոգիական հատկություններ կկառուցվեն հաջորդ թեմաներում:Իսկ մենք այստեղ սահմանափակվենք նշելով, որ այնպիսի երկրաչափական հատկություններ կամ մեծություններ ինչպիսիք են՝ անկյուն, հեռավորություն, բազմանկյուն, երկարություն, մակերես, ծավալ չեն հանդիսանում տոպոլոգիական հատկություններ կամ տոպոլոգիական ինվարիանտներ:



Not done:
Խնդրի շարունակում (Extension Problem)(Red text note)


:Տոպոլոգիական դասակարգում (Topological Classification)

s